\section{Alternatives and Justification}
\label{sec:alternatives}

\subsection{The Template Language}
\label{sec:language}

The course asks for Acceleo, and Acceleo exists in two incompatible
generations. Table~\ref{tab:languages} compares them with the two languages a
generator of this kind would otherwise be written in.

\begin{table}[!htbp]
\centering
\small
\begin{tabular}{>{\raggedright\arraybackslash}p{0.17\linewidth}>{\raggedright\arraybackslash}p{0.38\linewidth}>{\raggedright\arraybackslash}p{0.37\linewidth}}
\hline
\textbf{Language} & \textbf{For} & \textbf{Against} \\
\hline
Acceleo 3 (MTL) & Implements the OMG MOF Model-to-Text standard; the most documented generation. & The older generation, kept for compatibility; its generator project assumes an Eclipse workspace, and it is not the Acceleo of the release the build pins. \\
\textbf{Acceleo 4} & Expressions in AQL, the query language the Sirius tools share; a maintained standalone API; released in the same Eclipse release train as the Xtext, OCL and QVTo versions the build already pins. & A newer, thinner documentation; its whitespace rule (Section~\ref{sec:margins}) differs from every other template language. \\
Xtend & Typed template expressions in Java's type system. & Not a model-to-text language the course accepts; a template is Java code in disguise. \\
Java and a YAML library & The production implementation's own approach: build typed objects, let one serializer write the bytes. & Not model-to-text: a serializer owns the formatting, so the templates would decide nothing about the text, which is what the task asks for. \\
\hline
\end{tabular}
\caption{Template Languages Considered}
\label{tab:languages}
\end{table}

Acceleo~4 was chosen. The deciding fact is the release train: the build pins
one Eclipse release for every modelling tool (Task~1), and the Acceleo of that
release is Acceleo~4. Its standalone API runs the templates from a JUnit test or the
command-line pipeline with no workspace. Its expression language, AQL, is
close enough to OCL that the queries read like the constraints of Task~1.

\subsection{The Structure of the Templates}
\label{sec:structure}

Table~\ref{tab:structure} summarises the four decisions about how the templates
are organised. The subsections give the reasons.

\begin{table}[!htbp]
\centering
\small
\begin{tabular}{>{\raggedright\arraybackslash}p{0.2\linewidth}>{\raggedright\arraybackslash}p{0.33\linewidth}>{\raggedright\arraybackslash}p{0.37\linewidth}}
\hline
\textbf{Decision} & \textbf{Chosen} & \textbf{Rejected} \\
\hline
Who formats the text & The templates, character for character & A YAML serializer after the templates \\
Modules & One module, one template per kind of file, shared fragments & One module per kind of file \\
Files that several projects share & One query over every model rendered together & Merging the models in Java first, or a second transformation \\
Where spelling rules live & AQL queries in the module & Attributes the transformation computes \\
\hline
\end{tabular}
\caption{Structural Decisions}
\label{tab:structure}
\end{table}

\subsubsection{The Templates Write Every Character}
\label{sec:nopostprocessing}

The production implementation builds typed objects and lets one YAML
serializer write them, so its key order, indentation and quoting are the
serializer's. The Acceleo implementation could have generated a rough text and
run the same kind of serializer over it. That was rejected for two reasons.
First, it would move the formatting out of the model-to-text transformation,
which is what this task is meant to show. Second, it would weaken the
comparison: a serializer that normalises the output would make two different
texts equal, and the oracle would no longer prove that the templates are
right. The templates therefore write every character, and the test compares
their output with the oracle with nothing in between.

\subsubsection{One Module}
\label{sec:onemodule}

Every Kubernetes object carries the same metadata block and the same fixed set
of five labels, and several objects share larger fragments: the label set, a
network peer, a probe target, an environment list. With one module per kind of
file, each fragment would be written several times or imported across
modules. With one module, each fragment is one template, and the sixteen kinds
of file each have one template that spells their format. The module is long
(831 lines), but it is read top to bottom like the list of generated files.

\subsubsection{Files Several Projects Share}
\label{sec:shared}

Each entry point of the web proxy has one directory, and its index file lists
the routes of every project the entry point serves. The transformation resolves
one project at a time (Task~2), so each resolved model only knows its own
routes. Three ways to write the shared index were considered. A Java step could
merge the index entries of all models before the templates run, but that is a
post-processing step between the transformation and the templates. A second
transformation could combine the resolved models into one model of the whole
cluster, but that adds a metamodel for a single file. The chosen way is a query:
when several resolved models are rendered together, the index template lists
the routes that every model puts under the same path. Each model then writes
the same complete file, so the order in which the models are rendered does not
matter.

\subsubsection{Spelling Rules as Queries}
\label{sec:spelling}

Some values are spelled differently in YAML than in the model: a duration
\texttt{60s} is written as the integer \texttt{60}, a precedence of
\texttt{1} is written as the proxy priority \texttt{999}, and a string such as
\texttt{"50"} must be quoted, or a YAML reader takes it for a number. The
transformation could compute these spellings as extra attributes, but then the
resolved model would contain YAML details, which the target metamodel keeps
out on purpose. They are AQL queries in the module instead, one per rule, so
every rule about the text lives in the code generation.
